\section{Why climate-policy DID is especially hard to audit}
\label{app:climate}

Climate-policy DID studies concentrate several of the identification threats the
methodological literature already flags as fragile, which is what makes an
evidence-adequacy audit especially consequential here. \emph{Staggered and phased
rollouts} across jurisdictions invite the heterogeneous-timing failures of two-way
fixed effects. \emph{Spatial spillovers and cross-border leakage} under emissions
trading violate no-interference (SUTVA) when emitting or abating activity relocates
to control regions. \emph{Overlapping national and subnational regulations} are
routinely co-timed with the policy of interest, threatening the
no-confounding-policy assumption. \emph{Anticipatory abatement} in the gap between
a policy's announcement and its enforcement breaks no-anticipation. And
\emph{heterogeneous, noisy emissions measurement} (reporting-standard changes,
coverage differences, and estimation error) threatens consistent measurement of
the outcome. Each of these makes an identification assumption both more likely to
fail and harder to verify from a paper's reported evidence, which is why \sys{}
treats abstention (an explicit \texttt{unknown}) as a first-class output rather
than committing a confident but ungrounded judgement.

\section{Supplementary Figures and Tables}
\label{app:supp}

This appendix collects supporting material referenced from the main text:
the per-dimension risk distribution on the 27-paper real corpus
(Figure~\ref{fig:realcorpus}), and the deterministic calibration results against
the adjudicated human gold (Table~\ref{tab:calib}, Figure~\ref{fig:calib_tradeoff}),
and the agreement metrics for \sys{} versus the human gold (Table~\ref{tab:vsgold}).

\begin{table}[t]
  \centering\small
  \caption{\sys{}-LLM vs adjudicated human gold (55 cells, 5 papers) under two
  \texttt{unknown} policies: coverage (drop abstentions) and strict
  (\texttt{unknown} counts as mismatch). Pilot-scale: only 2 gold cells are
  \emph{high}; the robust signal is the one-directional over-severity and the
  abstention rate.}
  \label{tab:vsgold}
  \setlength{\tabcolsep}{4pt}
  \begin{tabular}{lcccc}
    \toprule
    \texttt{unknown} policy & $n$ & exact & $\kappa$ & wt.\ $\kappa$ \\
    \midrule
    exclude (coverage) & 33 & 0.152 & 0.019 & 0.044 \\
    mismatch (strict)  & 55 & 0.091 & 0.002 & -- \\
    \bottomrule
  \end{tabular}
\end{table}

\begin{figure}[t]
  \centering
  \includegraphics[width=\columnwidth]{figures/figure_vsgold}
  \caption{\sys{}-LLM versus adjudicated human gold on 55 dimension-level cells
  (5 papers). \textbf{(a)}~coverage, separating answered judgments from
  abstentions; \textbf{(b)}~severity distributions of human gold and \sys{};
  \textbf{(c)}~error direction among the 33 answered cells; \textbf{(d)}~the
  \sys{}-by-gold confusion matrix. The dominant signal is one-directional
  over-severity: \sys{} assigns higher risk than gold in 28/33 answered cases and
  lower risk in none.}
  \label{fig:vsgold}
\end{figure}

\begin{figure}[t]
  \centering
  \includegraphics[width=\columnwidth]{figures/figure_realcorpus}
  \caption{Per-dimension risk distribution on 27 real DID papers. The keyword
  baseline (left) is almost entirely \emph{low}; the
  two-stage LLM (right) is differentiated, with a large \texttt{unknown} share
  that marks retrieval failure rather than over-flagging.}
  \label{fig:realcorpus}
\end{figure}

\begin{table}[t]
  \centering\small
  \caption{Deterministic calibration vs the human gold (55 cells), two policies
  for weak-retrieval \emph{high}: demote to \emph{medium} (keep coverage) or
  abstain to \texttt{unknown}. Both cut over-severity and preserve high recall.
  Pilot scale: only 2 gold cells are \emph{high}.}
  \label{tab:calib}
  \setlength{\tabcolsep}{4pt}
  \begin{tabular}{lccc}
    \toprule
    metric & before & demote$\to$med & abstain$\to$unk \\
    \midrule
    over-severe (answ.)     & 28 & 18 & 6 \\
    high prec.\ / rec.      & 0.04 / 0.50 & 0.50 / 0.50 & 0.50 / 0.50 \\
    exact agr.\ (answ.)     & 0.15 & 0.45 & 0.54 \\
    wt.\ $\kappa$ (answ.)   & 0.06 & 0.13 & 0.25 \\
    answered / \texttt{unk} & 33 / 22 & 33 / 22 & 13 / 42 \\
    \bottomrule
  \end{tabular}
\end{table}

\begin{figure}[t]
  \centering
  \includegraphics[width=\columnwidth]{figures/calibration_tradeoff}
  \caption{\textbf{Calibration trade-off: severity reduction versus coverage}
  against the adjudicated human gold. Demoting weak-retrieval \emph{high}s to
  \emph{medium} preserves coverage (33 answered) while cutting over-severe
  judgments from 29 to 18; abstaining gives the strongest correction
  (over-severe 6, weighted $\kappa$ $0.25$) but leaves only 13 answered cells.
  Calibration corrects severity bias; the remaining bottleneck is evidence
  grounding, not severity thresholding alone.}
  \label{fig:calib_tradeoff}
\end{figure}
